\documentclass[12pt]{article}
\usepackage[utf8]{inputenc}
\usepackage[T1]{fontenc}
\usepackage{hyperref}
\title{Kadar Language Resources for K-12 Curriculum Development}
\author{Compiled by Claude Sonnet 5.5 (claude-sonnet-5-5) for Andr\'as Kornai}
\date{}
\begin{document}
\maketitle

\section{Introduction}
Kadar (ISO 639-3 kej) is a Dravidian language spoken by the Kadar people in Kerala and Tamil Nadu, close to Malayalam. Some scholars treat it as a separate language and others as a Malayalam dialect. One source gives 2,000 native speakers (2004 survey), and speakers are shifting to Tamil or Malayalam. The search returned almost entirely general-interest and encyclopedic pages. No Kadar-language teaching, dictionary or audio resources were found.

\section{Primary Educational Resources}
None found.

\section{Linguistic and Reference Materials}
\begin{itemize}
\item Language overview: \url{https://en.wikipedia.org/wiki/Kadar_language}
\end{itemize}

\section{Religious and Community Resources}
\begin{itemize}
\item Kadar people (Wikipedia): \url{https://en.wikipedia.org/wiki/Kadar_people}
\item Britannica entry on the Kadar: \url{https://www.britannica.com/topic/Kadar}
\end{itemize}

\section{Digital and Audio Resources}
None found.

\section{Curriculum Development Context}
Nothing found suggests any existing Kadar-medium curriculum. Starting one would require new documentation work with Kadar communities, ideally using Malayalam or Tamil school contexts, since Kadar children live within those systems.
\end{document}